\section{Model weights: representation, loading and use}
\label{sec:model-weights}

\subsection{Parameters and activations}

\textbf{Model weights} are learned numerical tensors: arrays with shapes
and element types that specify the model's transformations.  They
include embedding tables, attention projection matrices, feed-forward
matrices, learned normalization coefficients, and the output projection.
Some architectures tie the embedding and output tensors, so two uses need
not imply two independent allocations.  Ordinary inference uses the
trained parameters without updating them; no training gradients or
optimizer state are required for that computation.

A linear operation makes the distinction precise:
\begin{equation}
\label{eq:weights-activations}
 Y=XW,\qquad
 X\in\mathbb{R}^{m\times d_{\mathrm{in}}},\quad
 W\in\mathbb{R}^{d_{\mathrm{in}}\times d_{\mathrm{out}}},\quad
 Y\in\mathbb{R}^{m\times d_{\mathrm{out}}}.
\end{equation}
\eqcaption{A learned matrix \(W\) transforms \(m\) input activation
vectors into \(m\) new activation vectors.  Biases are omitted.}

Here \(W\) is persistent model data, whereas \(X\) and \(Y\) depend on
the positions being evaluated.  The same \(W\) can serve many positions
and many requests.  In attention, \(W_K\) and \(W_V\) are weights;
\(K=XW_K\) and \(V=XW_V\) are computed values.  Retaining these values
as KV state does not create additional trained parameters.  This
separation explains why the weight payload is principally a property of
the model and its representation, while KV payload also depends on the
workload's histories.

\subsection{From model files to runtime tensors}

A checkpoint or inference model file stores serialized tensors and
metadata describing such properties as names, shapes and numerical
representations.  Architecture configuration and tokenizer information
may be provided alongside it.  Large models can be split across several
files.  Files are persistent storage artifacts; kernels ultimately need
addressable tensor representations in the process's memory space.
\figref{fig:weight-loading} separates this preparation from repeated
execution.

One loading path reads tensor data and copies it into allocated buffers.
Another maps file-backed storage into the virtual address space.
\textbf{Memory mapping} establishes an addressable range without
necessarily populating all its physical pages immediately.  Accesses
can fault pages in on demand; prefetching and explicit population can
change when that cost occurs.  The runtime may use the mapped
representation directly or create additional tensors.  PyTorch's
checkpoint-loading documentation illustrates both mapped loading and
copying checkpoint tensors into a model's parameter allocations; these
are alternatives, not a universal loading sequence~\cite{pytorchloading2026}.

\begin{figure}[htbp]
\centering
% Vector schematic: preparation above, memory-to-cache-to-core traffic below.
% Icons are drawn locally so the figure stays scalable and matches the thesis.
\begingroup
\resizebox{\textwidth}{!}{%
\begin{tikzpicture}[
 font=\small\sffamily,text=black!80,
 data/.style={draw=figcol,fill=figcol!8,line width=0.65pt},
 mechanism/.style={draw=tabline,fill=neutralcol,line width=0.65pt},
 flow/.style={-{Stealth[length=5pt,width=4pt]},draw=black!65,line width=0.85pt},
 note/.style={font=\footnotesize\sffamily,align=center},
 label/.style={align=center,text width=2.9cm},
 heading/.style={font=\small\sffamily\bfseries,anchor=west}
]
\node[heading] at (-0.9,1.6) {1\quad Prepare addressable weights};

% Serialized tensor file: a folded document with a numerical array and metadata.
\begin{scope}[shift={(0.5,0.25)}]
 \draw[data] (-0.62,-0.72)--(0.62,-0.72)--(0.62,0.42)
   --(0.32,0.72)--(-0.62,0.72)--cycle;
 \draw[draw=figcol,line width=0.65pt] (0.32,0.72)--(0.32,0.42)--(0.62,0.42);
 \foreach \x in {0,1,2} \foreach \y in {0,1}
  \draw[draw=figcol!65,fill=white] (-0.4+\x*0.25,0.02+\y*0.23)
   rectangle ++(0.2,0.18);
 \draw[draw=figcol!65] (-0.4,-0.22)--(0.37,-0.22);
 \draw[draw=figcol!65] (-0.4,-0.42)--(0.12,-0.42);
\end{scope}
\node[label] at (0.5,-0.85) {Model files};
\draw[flow] (1.4,0.25)--(2.9,0.25);

% Loading makes tensor bytes addressable: storage and memory buffer symbols.
\begin{scope}[shift={(4,0.25)}]
 \draw[mechanism,rounded corners=1pt] (-0.9,-0.5) rectangle (-0.3,0.5);
 \foreach \y in {-0.3,0,0.3}
  \draw[draw=black!45] (-0.76,\y)--(-0.44,\y);
 \draw[flow] (-0.15,0)--(0.35,0);
 \draw[mechanism,rounded corners=1pt] (0.5,-0.5) rectangle (1,0.5);
 \foreach \y in {-0.25,0,0.25}
  \draw[draw=black!45] (0.5,\y)--(1,\y);
\end{scope}
\node[label] at (4,-0.85) {Read or map};
\draw[flow] (5.25,0.25)--(6.5,0.25);

% Optional preparation: rearrange a matrix into a kernel-friendly layout.
\begin{scope}[shift={(7.8,0.25)}]
 \foreach \x in {0,1,2} \foreach \y in {0,1,2}
  \draw[mechanism] (-1+\x*0.22,-0.34+\y*0.22) rectangle ++(0.18,0.18);
 \draw[flow] (-0.18,0)--(0.3,0);
 \foreach \x in {0,1,2,3} \foreach \y in {0,1}
  \draw[mechanism] (0.43+\x*0.19,-0.24+\y*0.28) rectangle ++(0.15,0.22);
\end{scope}
\node[label] at (7.8,-0.85) {Convert / pack};
\node[note,text=black!60] at (7.8,-1.35) {optional};
\draw[flow] (9.1,0.25)--(10.5,0.25);

% Runtime representation: overlapping tensor planes, not physical page residency.
\begin{scope}[shift={(11.7,0.25)}]
 \draw[data] (-0.42,-0.28) rectangle (0.88,0.72);
 \draw[data] (-0.65,-0.51) rectangle (0.65,0.49);
 \foreach \x in {1,2,3}
  \draw[draw=figcol!55] (-0.65+\x*0.325,-0.51)--++(0,1);
 \foreach \y in {1,2}
  \draw[draw=figcol!55] (-0.65,-0.51+\y/3)--++(1.3,0);
\end{scope}
\node[label] at (11.7,-0.85) {Runtime tensors};
\draw[draw=tabline,dashed] (-0.9,-1.9)--(13.2,-1.9);

\node[heading] at (-0.9,-2.6) {2\quad Fetch operands during inference};

% Main-memory pages: one contiguous row is selected for the illustrated fill.
\node[label,font=\small\sffamily\bfseries] at (0.5,-3.7) {Backing memory};
\begin{scope}[shift={(0.5,-5.15)}]
 \draw[data,fill=white] (-0.82,-0.78) rectangle (1.28,0.88);
 \draw[data,fill=white] (-1.0,-0.96) rectangle (1.1,0.7);
 \draw[data,fill=white] (-1.18,-1.14) rectangle (0.92,0.52);
 \foreach \x in {0,1,2,3} \foreach \y in {0,1,2,3}
  \draw[draw=figcol!35,fill=figcol!4]
   (-1.03+\x*0.48,-0.98+\y*0.35) rectangle ++(0.43,0.29);
 \foreach \x in {0,1,2,3}
  \draw[data,fill=figcol!30,line width=0.8pt]
   (-1.03+\x*0.48,-0.28) rectangle ++(0.43,0.29);
\end{scope}
\node[note] at (0.5,-6.65) {HBM and/or DDR};

% CPU boundary contains the cache hierarchy and the executing core.
\draw[draw=tabline,rounded corners=5pt,line width=0.8pt]
 (5.3,-3.85) rectangle (13.2,-7.1);
\node[note,anchor=west,font=\footnotesize\sffamily\bfseries]
 at (5.45,-3.6) {CPU (schematic)};

% Cache rows show retained copies; highlight only the illustrated line.
\begin{scope}[shift={(7,-5.15)}]
 \draw[mechanism,rounded corners=3pt] (-1.2,-1.05) rectangle (1.2,0.75);
 \foreach \x in {0,1,2,3} \foreach \y in {0,1,2}
  \draw[draw=tabline,fill=white]
   (-1.01+\x*0.51,-0.82+\y*0.46) rectangle ++(0.45,0.32);
 \foreach \x in {0,1,2,3}
  \draw[data,fill=figcol!30,line width=0.8pt]
   (-1.01+\x*0.51,-0.28) rectangle ++(0.45,0.32);
\end{scope}
\node[note] at (7,-6.65) {CPU caches};

% A cache-line packet on the data path makes the transfer granularity visible.
\draw[flow] (1.8,-5.27)--(5.78,-5.27);
\foreach \x in {0,1,2,3}
 \draw[data,fill=figcol!30,line width=0.8pt]
  (2.9+\x*0.44,-5.43) rectangle ++(0.39,0.32);
\node[note] at (3.75,-4.65) {Cache-line fill};
\node[note,text=black!60] at (3.75,-5.9) {on a miss to memory};

% Core package with pins and a small matrix multiplication kernel.
\begin{scope}[shift={(11.4,-5.15)}]
 \foreach \p in {-0.55,-0.18,0.18,0.55} {
  \draw[draw=black!50,line width=1pt] (\p,-0.9)--(\p,-1.08);
  \draw[draw=black!50,line width=1pt] (\p,0.9)--(\p,1.08);
  \draw[draw=black!50,line width=1pt] (-0.9,\p)--(-1.08,\p);
  \draw[draw=black!50,line width=1pt] (0.9,\p)--(1.08,\p);
 }
 \draw[mechanism,rounded corners=2pt] (-0.9,-0.9) rectangle (0.9,0.9);
 \draw[draw=black!30,fill=white] (-0.72,-0.72) rectangle (0.72,0.72);
 \foreach \x in {0,1} \foreach \y in {0,1} {
  \draw[mechanism] (-0.5+\x*0.19,0.06+\y*0.19) rectangle ++(0.15,0.15);
  \draw[data,fill=figcol!30] (0.16+\x*0.19,0.06+\y*0.19) rectangle ++(0.15,0.15);
 }
 \node[font=\small\sffamily] at (0,0.24) {$\times$};
 \node[note] at (0,-0.34) {Kernel};
\end{scope}
\node[note] at (11.4,-6.65) {CPU core};
\draw[flow] (8.22,-5.27)--(10.3,-5.27);
\node[note] at (9.3,-4.65) {Operands};

% Legend uses swatches, with an additional key for the highlighted data line.
\node[data,rounded corners=1pt,minimum width=0.28cm,
 minimum height=0.22cm,inner sep=0pt] at (-0.55,-7.85) {};
\node[note,anchor=west] at (-0.28,-7.85) {Weight data};
\node[mechanism,rounded corners=1pt,minimum width=0.28cm,
 minimum height=0.22cm,inner sep=0pt] at (3.6,-7.85) {};
\node[note,anchor=west] at (3.87,-7.85) {Preparation / hardware};
\node[data,fill=figcol!30,minimum width=0.28cm,
 minimum height=0.22cm,inner sep=0pt] at (9.3,-7.85) {};
\node[note,anchor=west] at (9.57,-7.85) {Selected weight data};
\end{tikzpicture}%
}
\endgroup

\caption{From model files to weight operands.  The upper path shows
preparation of addressable tensors: reading and mapping are alternatives,
and conversion or packing is optional.  Mapping does not imply immediate
physical residency.  Below, a highlighted row illustrates weight data
fetched from backing memory into CPU caches as a cache-line fill; a
kernel then uses its operands.  Cache hits reuse retained data without
another main-memory transfer.  Pages can become resident during startup
or on demand.  Array cells and hardware drawings are schematic; the cache
hierarchy and transfers within it are collapsed.}
\label{fig:weight-loading}
\end{figure}
\FloatBarrier

Representation preparation can convert element types, unpack or pack
quantized blocks, or arrange matrices in a layout suited to a kernel.
\textbf{Packing} changes the storage order or layout to make accesses
and arithmetic more efficient; it need not change the represented
values or their precision.  Such preparation can use temporary buffers
or retain an additional representation.  A loader's peak memory use can
therefore exceed the final tensor payload.  A short loading timer also
does not prove that every mapped page has already been read from storage.

The operating system backs virtual addresses with physical pages.
Where these pages reside depends on the allocation type, page-fault
behavior and applicable memory policy, as well as on the loading path.
File-backed mappings and anonymous tensor allocations need not obey the
same locality assumptions.  In particular, the thread that prepares a
representation and the threads that later read it may run on different
CPU locality domains.  Allocation locality is developed in
\secref{sec:cpu-memory-systems}.

\subsection{Repeated use and numerical representation}

Once prepared, weights can remain available across requests.  This
avoids repeating file loading, but it does not keep all weights
permanently in CPU caches.  Each kernel requests its operands through
the cache hierarchy; cache misses fetch data from the backing memory.
A large weight working set can require repeated main-memory traffic
across layer evaluations and generation steps.  Model loading is thus
a storage-to-memory preparation activity, whereas operand fetching is
part of ongoing computation.

Access behavior also depends on the operation.  A dense matrix product
uses its matrix entries across the evaluated activation rows.  An
embedding lookup selects rows for particular token IDs and does not
necessarily read the whole embedding table.  Sparse or conditionally
activated architectures can use subsets of their parameters.  It is
therefore unsafe to equate the full model size with the exact bytes read
on every generation step.

FP32 stores four bytes per scalar; FP16 and BF16 store two, with
different distributions of exponent and significand bits.  Quantization
represents values with fewer bits and uses associated information, such
as scale factors and sometimes zero points, to interpret groups of
values.  A nominal four-bit weight payload is half a byte per parameter,
but its format also stores metadata and may add padding.  Not all
tensors must use the same type.  Kernels may compute directly on a
compressed representation or reconstruct values as part of execution;
the accumulation type can differ from the stored type.  Storage savings
therefore do not by themselves establish either arithmetic speed or
numerical quality~\cite{llamacppproject2026}.

A parameter count describes how many learned scalar values a model has,
not its file size, process memory footprint, or inference traffic.  The
next sections identify the other data retained by inference, before
\secref{sec:memory-requirements} combines their capacity estimates.
\FloatBarrier
